\section*{Acknowledgements}
\addcontentsline{toc}{section}{Acknowledgements}

This proposal would not have come about without the prior work of
many people, whom it would be impractical to name exhaustively here
although most of their names occur \emph{passim} in the References
section of this document.

The Berlin/UCLA team led by Robert Englund and Peter Damerow, which
itself gave rise to CDLI \cite{CDLI}, most recently managed by Émilie
Pagé-Perron and Jacob Dahl and others, carried out foundational work
on the print and digital resources on which all work on
Proto-Cuneiform is based.

Michael Everson, Laura Hawkins and Debbie Anderson carried out and
fostered work on Proto-Cuneiform for Unicode under the auspices of
the Script Encoding Initiative
\cite{L2/16-267,L2/17-157,L2/19-284}. More recently, Anshuman Pandey
laid the basis for the current proposal in several important
proposal documents of his own \cite{L2/20-193,L2/22-239,L2/23-190}
and Robin Leroy was the principal author of a companion proposal to
this one, Archaic Cuneiform Numerals (ACN) \cite{L2/24-210}.

Steve Tinney authored the bulk of the current proposal and created the
Oracc PCSL project \cite{PCSL}, the digital corpus (based on CDLI) and
resources from which the appendices are derived.

Robin Leroy authored the section \hyperref[rationale]{``Rationale for
  Separately Encoding Proto-Cuneiform''} and frequently shared his
expertise on cuneiform, Unicode, and LaTeX.

The PCSL font was initially created by Anshuman Pandey based on Bob
Englund's collection of signs and was an essential basis for the
current proposal. The font was revised and augmented by Steve
Tinney.
